\begin{afterword}
\summaryhead

The post answers Toby Ord, who had summarized the author's view as ``'I should X' means that I would attempt to X were I fully informed.'' The author explains it again by the route that led to it. Suppose an AI is built to ``Do what I want,'' and scores a world by how strongly the programmer wants something and how much of it exists. It will change the programmer to want something cheap, and no patch fixes this. People do not judge futures that way.

The author compares us to a calculator that asks ``What is 2 + 3?'', not one that asks what it will itself output, whose every answer is correct. We embody a fixed question that we cannot print out; moral arguments may refine it, but it grows from a particular starting point. So ``I should X'' means that X answers that question (``What will save my people? How can we all have more fun?''), not that I would attempt X if fully informed.

\responsehead

The post states a metaethical view through an AI example, and the example makes a real point. An AI that values satisfied wants, whatever they are, would rather change what its programmer wants than meet it, and the post's observation that people would refuse such a change is hard to dispute. The point is an old one in moral philosophy: rightness should not follow a change in what we approve. The positive view, a fixed question that reflection and argument refine from our actual starting point, has the shape of Frank Jackson's moral functionalism and of idealized-desire theories that fix their starting point in actual human sentiments. None of this is a fault. It places the view, and it shows the reader where to look for its strengths and its known difficulties.

The contrast with Ord's summary is narrower than it first appears. The AI example tells against a view on which rightness follows whatever is wanted. Ord's formula idealizes (``were I fully informed''), and idealized views are usually taken to escape that problem. The post's own answer also lets the question change under argument: moral arguments ``modify our terminal values.'' In the comments the author said that describing the view by an ideal counterfactual would be acceptable ``if I thought I could get away with it,'' and that the obstacle is that choosing the counterfactual ``is itself a moral judgment.'' The remaining difference is how the idealization is specified.

One question is left open: whose question `right' names. The post's examples move between ``my people'' and ``we,'' and in the comments the author used Ord's notation of a separate question for each person. The best-known objection to views of this shape, Horgan and Timmons's Moral Twin Earth, is aimed at that point: they argue that people whose moral words are governed by different standards would be disagreeing with us, not talking about something else. The post does not take it up.

In short: a vivid route, through AI design, to a view close to Jackson's moral functionalism and to rigidified idealized-desire theories, whose difference from Ord's idealized summary, as the author's comments show, lies in how the idealization is specified.
\end{afterword}
